\section{Results} \label{Res}

In this section, we show the results of the analysis of PR4 intensity and E-mode polarization data. For each case, we present the p-values and directions. We also re-examine the HPA in the PR3 data set using our alternative inpainting approach and to check consistency with previous results.

\subsection{Intensity results for PR4}

Following previous temperature analysis, we present the results of full-resolution, $N_{side} = 2048$, PR4 data for a set of discs with radii between 4 and 40 degrees. The p-values are defined as the fraction of simulations with local-variance dipole amplitude larger than the one observed in the data. Unlike PR3, where cleaned maps were provided for four methods, for PR4, CMB maps are only available for Sevem and Commander. As expected, both of them are in good agreement with the Planck 2018 analysis. For Sevem, none of the 600 simulations have an amplitude as large as the one observed in the data for discs of 4, 6, and 8 degrees. This means that, under the $\Lambda$CDM model assumption, the probability of having such asymmetry in the temperature sky in that range of angular scales is below 0.17\%. For radii larger than $8^{\circ}$ the p-value increases systematically. We have checked with simulations that have a dipolar modulation pattern, with A = 7\% and p = ($209^{\circ}, -15^{\circ}$), that this increase of the p-value with the radius of the disc is expected. We obtain the same behavior for Commander, where none of the 100 simulations shows an amplitude as large as the one measured in the data. {Note that due to the lower number of simulations available for Commander, our sensitivity for the p-value is 1\% for this method.

Taking into account that the asymmetry only appears at large angular scales, we repeat the analysis for low resolution maps, at $N_{side} = 64$. Figure \ref{fig:T-results} shows the p-value as a function of the disc radius for both component separation methods and both resolutions. For low resolution maps the p-value is computed using discs with radii between 4 and 90 degrees as in \cite{Planck:2019evm}. In general, the significance level is lower than in the full-resolution analysis. Nevertheless, for discs of $4^{\circ}$ and $6^{\circ}$, we continue to find that no simulation has an amplitude equal or greater than that of the data.

\begin{figure}[t!]
  \centering
  \includegraphics[scale=0.4]{Figuras/P_values_Temperature.pdf}
  \caption{p-values for the asymmetry measured through the local-variance estimator for PR4 Sevem and Commander temperature maps and the two considered resolutions, $N_{side}$ = 2048 and $N_{side}$ = 64. The p-value is inferred by checking how many simulations have a local dipole amplitude larger than the one observed in the data. The sensitivity depends directly on the number of available simulations, corresponding to 0.17\% for Sevem and 1\% for Commander.}
  \label{fig:T-results}
\end{figure}

Regarding the direction of the dipole, in Table \ref{tab:dirT} we provide the full-resolution and low-resolution direction for Sevem and Commander obtained by considering $4^{\circ}$ discs. In Sevem the observed direction is in excellent agreement with the previously reported one (see Table 22 from \cite{Planck:2019evm}), while in Commander there is a small discrepancy of a few degrees, which could be due to the difference in the number of simulations between releases. In particular, we have checked that the $\ell$ and b angles have a very large dispersion (see Fig. \ref{AngularDistance_Sensitivity}), especially when the amplitude of the modulation is small, making consistent the values for Commander between both releases.

\begin{table}[t!]
\centering
\begin{tabular}{ccccc}
\cline{2-5}
          & \multicolumn{4}{c}{($\ell$, b) {[}deg{]}}                                     \\ \cline{2-5} 
          & \multicolumn{2}{c}{PR3}               & \multicolumn{2}{c}{PR4}               \\ \hline
Data      & $N_{side}$ =  64 & $N_{side}$ =  2048 & $N_{side}$ =  64 & $N_{side}$ =  2048 \\ \hline
Sevem     & (208, -15)       & (205, -20)         & (209, -14)       & (205, -19)         \\
Commander & (213, -16)       & (207, -20)         & (209, -15)       & (205, -20)         \\ \hline
\end{tabular}
\caption{Local-variance dipole directions for PR4 Sevem and Commander temperature maps. Directions are measured in galactic coordinates for both resolutions and using $4^{\circ}$ discs.}
\label{tab:dirT}
\end{table}

\subsection{E-mode polarization results for PR3}

In order to compare the performance of our GCR inpainting technique versus the purified inpainting used in the Planck 2018 analysis \cite{Planck:2019evm}, we have first carried out the analysis of PR3 data using our method.
We have followed the procedure explained in Section \ref{subsection:3.2}. We recall that, in this case, each component separation algorithm has in total 300 noise E2E simulations each of them accompanied by 3 independent CMB skies. We have done 30 splits of two independents sets of 150 noise simulations and computed the covariance matrices as previously explained in section \ref{subsection:3.2}. Combining the 150 noise simulations with the 450 independent CMB skies, for the analysis we have available 60 sets of 450 simulations. 

Table \ref{PR3_Results} summarizes the results of this work. For each component-separation algorithm we provide the preferred direction of the local variance dipole and the p-values obtained using the reference mask. In addition, in Table \ref{Alignment_PR3} we give the angular distance between the temperature and E-mode dipoles, and the p-value associated with this alignment. 
As one would expect, we recover results quite similar to those obtained by the Planck collaboration, validating also our approach. We can distinguish two groups. While for Sevem and Commander the p-value of having such a large amplitude is below 1\%, for SMICA and NILC the probability is at the level of 3-4\%. Regarding the TE alignment, we are below 6\% for all the component separation methods. 
Again results are overall consistent between both analysis, although the result for NILC and SMICA becomes less anomalous in our case (p-values of 5.7 and 2.0 versus 1.9 and 0.9, respectively). These small deviations can be explained by the fact that our analysis presents some differences with respect to the previous one, such as the inpainting technique, a different mask or the use of splits.

\begin{table}[t!]
\centering
\begin{tabular}{ccccc}
\hline
Data      & p-value {[}\%{]} & p-value range {[}\%{]}       & ($\ell$, b) {[}deg{]} & p-value {[}\%{]} \\
& (this work) & & &  (from \cite{Planck:2019evm} )
\\ \hline
Sevem     & 0.22                         & {[}\textless{}1/450, 0.44{]} & (232, -9) $\pm 4$     & 0.4              \\
Commander & 0.70                         & {[}\textless{}1/450, 1.2{]}  & (222, -9) $\pm 4$     & 0.7              \\
SMICA     & 4.4                          & {[}1.8, 6.4{]}               & (225, -12) $\pm 4$    & 5.5              \\
NILC      & 3.4                          & {[}1.6, 5.0{]}               & (238, -16) $\pm 5$    & 5.8              \\ \hline
\end{tabular}
\caption{Local-variance dipole directions and p-values for the PR3 four component-separated E-mode polarization maps, analysed with the inpainting procedure, at $N_{side}$ = 64, together with the 1$\sigma$ interval obtained from the 60 sets as explained in Section \ref{subsection:3.2}. Measured directions are also showed. The error is estimated from the distribution of the angular distances between the mean direction and the directions of the 60 data sets, taking into account the contour, with azimuthal symmetry, that includes 68\% of the directions. All the values have been obtained using $4^{\circ}$ discs. Last column shows the p-value obtained in the previous work by the Planck collaboration \cite{Planck:2019evm}.}
\label{PR3_Results}
\end{table}

In order to test the robustness of the p-value against the mask, we have generated 4 additional masks using thresholds in the maximum reconstruction error (see Section \ref{subsection:3.2}) ranging from 45\% to 35\% and with the following allowed fractions of sky: 64.5\%, 63\%, 59.6\%, 57.6\%. Figure \ref{P_values_PR3_Against_Mask} presents the p-values obtained for the four component-separated E-mode polarization maps as a function of the $f_{\rm sky}$. The minimum p-value is obtained for the mask with $f_{\rm sky}$ = 59.6\%. Sevem and Commander are still below 1\% and the p-value is even lower than with the reference mask ($f_{sky} = 61.4\%$). However, SMICA and NILC present still values around or above a few per cent for all the considered masks. Additionally, Figure \ref{Dir_PR3_Against_Mask} displays the coordinates of the dipole direction as a function of the $f_{\rm sky}$. All the methods and masks show certain level of robutness in the sense that all the directions are concentrated in a small region near ($\ell$, b) = ($230^{\circ}$, $-13^{\circ}$). Furthermore, it seems that the latitude (b) is somehow more robust that the longitude ($\ell$).

\begin{table}[t!]
\centering
\begin{tabular}{ccccc}
%\cline{2-5}
\hline
          & \multicolumn{2}{c}{TE alignment (this work)} & \multicolumn{2}{c}{TE alignment (from \cite{Planck:2019evm})} \\
%          & \multicolumn{2}{c}{Alignment with T (this work)} & \multicolumn{2}{c}{Alignment with T (previous work)} \\\hline
Data      & $\cos{\alpha}$      & p-value {[}\%{]}     & $\cos{\alpha}$      & p-value {[}\%{]}     \\ \hline
Sevem     & 0.91                & 4.2                  & 0.86                & 6.9                  \\
Commander & 0.99                & 0.7                  & 0.99                & 0.9                  \\
SMICA     & 0.96                & 2.0                  & 0.99                & 0.9                  \\
NILC      & 0.88                & 5.7                  & 0.97                & 1.9                  \\ \hline
\end{tabular}
\caption{Angular distance of alignment between the preferred direction of the local variance dipole observed in intensity and the one for the E-mode polarization data. We use the mean directions presented in Table \ref{PR3_Results} to compute the angular distance, and the p-value is obtained using the E2E simulations combining 
all the results from different splits in a single distribution. 
The previous results by the Planck collaboration \cite{Planck:2019evm} are shown in the last column .}
\label{Alignment_PR3}
\end{table}

\begin{figure}[t!]
  \centering
  \includegraphics[scale=0.4]{Figuras/P_values_PR3.pdf}
  \caption{p-values for variance asymmetry as a function of $f_{\rm sky}$ for the four PR3 component-separated E-mode polarization maps analysed with the inpainting procedure.: SMICA (blue), Sevem (red), NILC (green) and Commander (orange). All values have been obtained at $N_{side} = 64$ and using $4^{\circ}$ discs. We recall that our reference mask for PR3 corresponds to a $f_{\rm sky}$=61.4\%.}
  \label{P_values_PR3_Against_Mask}
\end{figure}

\begin{figure}[t!]
    \centering
    \includegraphics[scale = 0.5]{Figuras/Direction_SMICA_SEVEM_NILC_COMMANDER_PR3.pdf}
    \caption{Local-variance dipole directions for the four PR3 component-separated E-mode polarization maps analysed with the inpainting procedure.: SMICA (circle), Sevem (cross), NILC (square) and Commander (triangle). Different colors corresponds to different $f_{sky}$: red (64.5\%), blue (63\%), green (61.4\%), orange (59.6\%), and black (57.6\%). All values have been obtained at $N_{side} = 64$ and using $4^{\circ}$ discs.}
    \label{Dir_PR3_Against_Mask}
\end{figure}

\subsection{E-mode polarization results for PR4}\label{PR4_E_Results}

Finally, we present the results for the Sevem PR4 data set using the GCR inpainting \footnote{Note that we have not used the Commander PR4 data since CMB and noise simulations are not provided separately
(i.e. only CMB+noise simulations are provided), which makes not possible to compute independently the noise covariance matrix. 
Although, in principle, it would be possible to estimate the total covariance matrix, this introduces additional uncertainties in the matrix elements. Moreover, only 400 simulations are available for polarization, increasing even further the error in the analysis.}.
We have followed the same procedure as for PR3 with a total of 
30 splits of 300/300 simulations.

Table \ref{P_values_PR4} summarizes the results, providing the preferred direction of the local variance dipole and p-values obtained for the Sevem algorithm using the PR4 reference mask. The error in the direction is obtained as the 68 per cent (single-tailed) of the distribution of the angular distance between the mean direction and the direction of the 60 data sets. Therefore, it is only given to provide an estimation of the dispersion in the different considered splits. In practice the error associated to the estimated direction should be significantly larger, since this quantity is below the errors inferred from the methodology itself or the effect of the inpainting. 
We also give in the same table the angular distance between the temperature and E-mode dipole directions, and the p-value associated with this alignment. The analysis shows that the p-value of the amplitude for Sevem PR4 increases with respect to the one obtained for PR3 (see Table \ref{PR3_Results}), increasing from 0.22 to 2.8 per cent. The p-value of the TE alignment is very similar for both pipelines and remains below 5 per cent.

\begin{table}[t!]
\centering
\begin{tabular}{ccccc}
%\cline{4-5}
\hline
%                 &                        &                       & \multicolumn{2}{c}{TE alignment}     \\ %\hline
p-value {[}\%{]} & p-value range {[}\%{]} & ($\ell$, b) {[}deg{]} & 
\multicolumn{2}{c}{TE alignment} \\ 
                &                        &                       &      $\cos{\alpha}$ & p-value {[}\%{]} \\ \hline
2.8              & {[}1.8, 3.8{]}         & (234, -14) $\pm$ 5    & 0.91           & 4.5              \\ \hline
\end{tabular}
\caption{Local variance dipole direction and p-value for the Sevem PR4 E-mode polarization map, analysed with the inpainting procedure, at $N_{\rm side} = 64$, together with the $1\sigma$ interval obtained from the 60 data sets.
Measured directions are also showed. The error is estimated from the distribution of the angular distance between the mean direction and the directions of the 60 data sets,
taking into account the contour, with azimuthal symmetry, that includes 68\% of the directions. All the values have been obtained using $4^{\circ}$ discs.}
\label{P_values_PR4}
\end{table}

Another interesting result is the level of the modulation in the data. Using the modulated simulations described in Section \ref{subsection:3.4}, the relation between the amplitude of the modulation (given by equation \ref{DM}) and the one measured in the local-variance map can be estimated. Figure \ref{Calibration} shows this relation together with the amplitude observed in the data local-variance map for the reference mask. 
According to this, for modulations with amplitudes larger than 6\%, the relation is linear, although with a large dispersion due to the E-mode noise level. Assuming that the modulation model (Eq.\ref{DM}) is the correct one, the asymmetry observed in the data corresponds to a modulation at the level of 9\%. Moreover, the estimated amplitude $A_{\rm LV}$ is within the 68\% of the distribution for the dipolar modulation amplitude ($A_{\rm DM}$) for the range between 6\% and 13\%.
%
\begin{figure}[t!]
    \centering
    \includegraphics[scale = 0.4]{Figuras/Rel_Amp.pdf}
    \caption{Relation between the dipolar modulation amplitude $A_{\rm DM}$ (input value in the modulated simulations) and the one measured in the local variance map $A_{\rm LV}$, assuming the realistic PR4 noise and systematics.
    For each amplitude, all the 60 inpainted data sets
    are used to generate a single distribution. For amplitudes larger than 6\%, there is a linear relation between both amplitudes with a slope of 0.046 (dashed black line). The horizontal black line corresponds to the mean amplitude (over the 60 data sets) measured in the data local-variance map using the PR4 reference mask. The measured $A_{\rm LV}$ value corresponds to
    a modulation at the level of 9\%, and it is within the 68\% of the distribution for the $A_{\rm DM}$ interval between 6\% and 13\%. Black dots in the tails of distributions represent the values outside the 99.7\%.}
    \label{Calibration}
\end{figure}

As before, we test the robustness of the p-value against the different masks. We have generated 6 additional masks using thresholds in the range from 45\% to 30\% with the following allowed fractions of sky: 67\%, 65.9\%, 63.2\%, 61.5\%, 59.4\%, and 57\% (we recall that the reference mask corresponds to f$_{\rm sky}$ = 64.6\%}).
Imposing more strict conditions in the error of the recovered E-map (i.e. lower thresholds) while keeping a reasonable $f_{\rm sky}$ is possible in this case because a more accurate E-mode reconstruction can be obtained in PR4. This is due to both, a lower level of systematics and twice the number of simulations that have been used for estimating the noise covariance matrix. Figure \ref{P_values_PR4_Against_Mask} presents the obtained p-values as a function of $f_{\rm sky}$. The minimum p-value (1\%) is obtained for the mask with $f_{\rm sky}$ = 61.5\%. For the 5 smallest masks, allowing a sky between 67\% and 61.5\%, the p-value is quite stable with a value below 2.8\%, and then it starts to rise with $f_{\rm sky}$. The largest mask returns a p-value of 7.4\%. 
This behaviour could be due to the loss of information as the mask increases, assuming the asymmetry had a cosmological origin.
However, it could also be due to the fact that a larger mask reduces foreground contamination, what we would expect if the asymmetry were due to the presence of residuals.
Additionally, Figure \ref{Dir_PR4_Against_Mask} and \ref{P_value_Alignment} display the directions of the local-variance dipoles and the p-value of the TE alignment, respectively, as a function of the $f_{\rm sky}$. 
In particular, the minimum p-value (4.6\%) is obtained for the reference mask. Finally, Figure \ref{Summary_Dir} summarizes in a Mollweide projection all the directions measured in this work.

\begin{figure}[t!]
  \centering
  \includegraphics[scale=0.4]{Figuras/P_values_PR4.pdf}
  \caption{p-values for variance asymmetry as a function of $f_{\rm sky}$ for the PR4 Sevem E-mode polarization map analysed with the inpainting procedure. The grey region corresponds to the 68\% interval obtained from the 60 data sets. All values have been obtained at $N_{\rm side} = 64$ and using $4^{\circ}$ discs.}
  \label{P_values_PR4_Against_Mask}
\end{figure}

\begin{figure}[t!]
    \centering
    \includegraphics[scale = 0.5]{Figuras/Direction_SEVEM_PR4.pdf}
    \caption{Local-variance dipole directions for the Sevem PR4 
    component-separated E-mode polarization maps analysed with the inpainting procedure. 
    Different colors corresponds to different f$_{\rm sky}$: red (67\%), blue (65.9\%), green (64.6\%), orange (63.2\%), black (61.5\%), pink (59.4\%), and brown (57\%). All values have been obtained at Nside = 64 and using 4$^\circ$ discs}
    \label{Dir_PR4_Against_Mask}
\end{figure}

\begin{figure}[t!]
    \centering
    \includegraphics[scale = 0.4]{Figuras/P_values_PR4_TE_Alignment.pdf}
    \caption{p-value of the alignment between the dipole axis in PR4 temperature and E-mode as a function of $f_{\rm sky}$.}
    \label{P_value_Alignment}
\end{figure}

\begin{figure}[t!]
    \centering
    \includegraphics[scale = 0.5]{Figuras/Direction_Summary.pdf}
    \caption{Local variance dipole directions for the PR3 four component-separated E-mode polarization maps analysed with the inpainting procedure: Sevem (red), SMICA (blue), NILC (green), and Commander (orange). The black dot corresponds to the direction measured using the Sevem PR4 polarization data. For reference, we also show the CMB dipole direction, and the directions observed in the temperature data for Sevem PR4 at $N_{side} = 64$ (red star) and $N_{side} = 2048$ (black star).}
    \label{Summary_Dir}
\end{figure}
